\begin{frame}[t]{Best Classical vs Best Deep}
  \sectiontag{Comparison}
  \takeaway{The final comparison is informative but not apples-to-apples: deep models benefit from adaptive splits, while classical models stay on the harsher fixed-cutoff setup.}
  \vspace{0.35em}

  \begin{columns}[T,onlytextwidth]
    \column{0.60\textwidth}
    \begin{panelbox}
      \textbf{Best notebook skill by representative series}
      \begin{center}
        \begin{tikzpicture}
          \begin{axis}[
            ybar,
            width=0.94\linewidth,
            height=0.58\textheight,
            ymin=0,
            ymax=0.9,
            ylabel={Weighted skill},
            symbolic x coords={High wt., Long hist., Stable, Volatile},
            xtick=data,
            bar width=12pt,
            enlarge x limits=0.18,
            ymajorgrids=true,
            grid style={draw=DeckBorder},
            axis line style={draw=DeckBorder},
            tick style={draw=DeckBorder},
            xticklabel style={font=\small},
            yticklabel style={font=\small},
            legend style={draw=none,fill=none,at={(0.5,1.06)},anchor=south,legend columns=-1,font=\small},
          ]
            \addplot[fill=DeckAccent!80,draw=DeckAccent!80] coordinates {
              (High wt.,0.0000)
              (Long hist.,0.0000)
              (Stable,0.0000)
              (Volatile,0.8412)
            };
            \addplot[fill=DeckMuted!55,draw=DeckMuted!55] coordinates {
              (High wt.,0.5117)
              (Long hist.,0.0348)
              (Stable,0.0000)
              (Volatile,0.8194)
            };
            \legend{Best classical, Best deep}
          \end{axis}
        \end{tikzpicture}
      \end{center}
    \end{panelbox}

    \column{0.36\textwidth}
    \begin{panelbox}
      \textbf{What the chart really says}
      \small
      \begin{itemize}
        \item deep helps most on the highest-weight case
        \item the volatile case stays close to classical ARIMA
        \item the stable case does not reward extra complexity
      \end{itemize}
      \divider
      \textbf{Caveat}
      Read this as a four-series notebook comparison under different split rules.
    \end{panelbox}
  \end{columns}
\end{frame}

\begin{frame}[t]{Final Conclusions}
  \sectiontag{Close}
  \takeaway{Across the three notebooks, the main lesson is that diagnostics and evaluation design matter as much as model family choice.}
  \vspace{0.4em}

  \begin{panelbox}
    \small
    \begin{enumerate}
      \item \textbf{Weighting and chronology define the task.} Any fair read must respect the temporal suffix split and the extreme concentration of row weights.
      \item \textbf{EDA narrows what is defensible.} Cutoff weakness, mixed raw-level stationarity, and short-memory lag structure all constrain the sensible notebook experiments.
      \item \textbf{Classical models stay credible.} Short histories favor simple smoothers, while the volatile series rewards differenced ARIMA.
      \item \textbf{Deep models help selectively.} Adaptive-split recurrent models improve some cases, but the split mismatch blocks a universal win claim.
    \end{enumerate}
  \end{panelbox}
  \vspace{0.15em}
  {\small\color{DeckMuted}\textbf{Closing note.} The notebooks contribute a disciplined path from diagnostics to honest interpretation, not a single universal winner.}
\end{frame}

\begin{frame}[plain]
  \vfill
  \begin{center}
    {\fontsize{28}{32}\selectfont\bfseries Questions?\par}
    \vspace{0.5em}
    {\large\color{DeckMuted}Thank you.}
  \end{center}
  \vfill
\end{frame}
